\scanpage{305}
\mspage{4}
de la forme
\[
 {{A^*}_*}^{**}=\bigl(( (A^*)_* )^*\bigr)^*
 \deleted{\unsure{ce sont}}\quad\text{p.ex.}
\]
--- dont les objets sont notés par des lettres à étoiles
$ {{K^{**}}_{**}}^*\ldots\ \Big]$

Le rôle de ${\Delta_{\unsure{*}}}^{*}\in\operatorname{Ob}\widehat{\boldsymbol\Delta}^{*}$
\inserted{$=\operatorname{Ob}Ss^*$}
comme ``objet cosimplicial universel'' est explicité par la

\uline{Proposition 1}. Soit $A$ une catégorie où les $\underline{\lim}$ existent
(resp. $B$ une catégorie où les $\varprojlim$ existent). Alors
\[
 \text{a)}\quad A^*\xleftarrow{\approx}
 \underline{\operatorname{Hom}}_{\text{comm. à }\varinjlim}(Ss,A)
\]
par $F\longmapsto F(\Delta^*)$
\[
 \text{b)}\quad B_*\xleftarrow{\sim}
 \underline{\operatorname{Hom}}_{\text{comm. à }\varprojlim}(Ss^o,B)
\]
par $F\longmapsto F((\Delta^*)^o)$

(où $(\Delta^*)^o$ désigne l'objet simplicial de $\widehat{\boldsymbol\Delta}^{o}$
défini par l'objet cosimplicial $\Delta^*$ de $\widehat{\boldsymbol\Delta}$).

Bien sûr, b) n'est autre que a) appliqué à $A=B^o$, et a) signifie que
le foncteur ``restriction à $\boldsymbol\Delta$''
\[
 \underline{\operatorname{Hom}}(\widehat{\boldsymbol\Delta},A)
 \longrightarrow\underline{\operatorname{Hom}}(\boldsymbol\Delta,A)
\]
induit une équivalence sur la ss-cat
\deleted{\unsure{des}}\inserted{comm. à $\underline{\lim}$}
du premier mb. C'est vrai pour \uline{ttes}
(\deleted{petites}\inserted{\unsure{nos}}) catégories $C$ ---
pas seulement pour $\boldsymbol\Delta$, et bien connu.
Le foncteur en sens inverse associe, à
$K_*=\varphi:\boldsymbol\Delta\mathbin{\text{---}}A$, le foncteur
\[
 X_*\longmapsto
 \underset{\Delta^n_*\in\boldsymbol\Delta/X_*}{\underline{\lim}}
 \varphi(\Delta^n_*)=
 \underset{\Delta^n_*\in\boldsymbol\Delta/X_*}{\underline{\lim}}K_{[n]}
\]
Dans le cas b) où on \unsure{écrit}
\[
 \underline{\operatorname{Hom}}(\widehat{\boldsymbol\Delta}^{o},B)
 \longrightarrow\underline{\operatorname{Hom}}(\boldsymbol\Delta^{o},B)
\]
le foncteur en sens inverse associe à
$K^*=\psi:\boldsymbol\Delta^o\mathbin{\text{---}}B$
le foncteur contravariant
\[
 X^*\longmapsto
 \varprojlim_{\Delta^n_*\in(\boldsymbol\Delta/X_*)^o}\varphi(\Delta^n)
 =\varprojlim_{\Delta^n\in(\boldsymbol\Delta/X_*)^o}K^{[n]}
\]

\uline{Yoga}: nous nous intéresserons surtout
\deleted{aux} (en première étape) aux invariants covariants attachés à $X_*$,
à valeurs dans une cat. $A$, qui \deleted{pour $X_*$}\inserted{commutent}
aux $\underline{\lim}$ p.r. à $X_*$, --- ils sont donc déterminés
(si $A$ est ``\unsure{grosse|grande}'' pour contenir les $\underline{\lim}$
(petites)) par des éléments de $A^*$ i.e. des objets \uline{cosimpliciaux} de $A$.
Pour les invariants contravariants en $X_*$, à val. dans un $B$, nous nous
attacherons en premier lieu à ceux qui \deleted{commutent} aux $\varprojlim$
en $X_*$; pour $B$ stable par $\varprojlim$, ils correspondent
\deleted{donc} aux objets de $\underline B_*$, i.e. les objets
$\tfrac12$ simpliciaux de $B$. Si $B=(Ens)$, cela correspond donc à des
ens. $\tfrac12$ simpliciaux, i.e. à des objets $K_*$ de
$\unsure{Ss|B}$ lui-même.
\margin{306}
